\chapter{Claim and Scope Boundaries}
\label{app:claims}

The project evaluates one pinned model, two simulators with specified hidden states and observation rules, and one historical educational dataset. It does
not provide evidence of improved human learning, mastery, motivation, reduced cognitive offloading, classroom readiness, or
realised deployment savings. The tutoring interface is a proof of concept. The exploratory CSEDM analysis measures prediction
and selection of likely errors. No corrective review intervention was performed.

\section{Limits of Each Study}

\paragraph{Benchmark of an external model (RQ1).}
The 23 eligible cases describe one fixed authored corpus and one recorded generation per case and channel. The evaluation model
is not a validated substitute for a student. The complete harness correction supersedes the original apparent advantage.
The remaining one-case accuracy difference leaves a useful benefit unresolved.
Passing an authored test also does not establish compliance with every task requirement. A favourable worked example cannot
establish that relevance caused the improvement, especially when unrelated passing evidence yields the same prediction.

\paragraph{Bounded tracking (RQ2).}
The positive finding concerns the declared average adverse-condition Brier score and clean-condition penalty rule. Missing and
contradictory evidence individually worsened. The bound limits one observation's contribution to log-odds. Repeated bounded
updates can accumulate. This is separate from an empirical accuracy or calibration claim.

\paragraph{Probe selection within a fixed budget (RQ3).}
The experiment selects a whole batch from calibrated starting probabilities. It does not test uncertain initial estimates,
conversational adaptation, or stopping when further evidence seems unhelpful. The candidate must fill its allocation.
Its clipped-belief score need not equal the expected error reduction of the decision executed, so comparisons apply to that
heuristic and the implemented baselines. The true-channel reference is not a proven optimum.

Average harm across the selected held-out settings does not mean every setting worsened. Their equal weights define the study
comparison, rather than the prevalence of failures in deployed tutors. Primary intervals condition on the fitted calibration
model. The per-environment point-estimate check is not a statistical safety guarantee. More simulated episodes improve precision
within these assumptions. They do not validate the assumptions against learners.

\paragraph{Finding errors in historical learner records (RQ4).}
The 729 targets come from 86 learners and 19 problems. Learner-separated folds support the specified prediction evaluation,
without establishing generalisation to unseen problems or courses. The uncertainty result concerns which existing model errors
are selected. It does not show that another observation or reviewer would correct them. Learner resampling retains fitted
predictions and selection identities. The interval excludes retraining and reranking variation, and does not fully represent
dependence from overlapping cross-validation training sets.

\section{Research and Demonstration}

The recorded experiment view explains frozen evidence. The optional live route, if retained, illustrates provider and sandbox
integration on a development case. It uses a separate simple scoring rule. A successful live run adds no canonical research
evidence, and neither route establishes tutoring efficacy. A fixed probe count measures allocation, not optimisation of varying
monetary or latency costs. Financial savings cannot be inferred from fewer simulated probes alone.

\section{Changes from the Proposal}

The original proposal included a broader multi-agent POMDP, Cognitive Bandwidth and Friction variables, sequential Monte Carlo,
reinforcement learning, and a future human pilot. These elements were narrowed or deferred because no validated learner transition
model or approved human-study design was available. The final technical work instead concentrates on measurable evidence handling,
controlled comparison, and reproducibility.
